\scanpage{81}
\mspage{40}
de type fini sur $Q$, $\operatorname{Pic}(S)$ est un $\mathbb Z$-\deleted{groupe}\inserted{module} de type fini \inserted{(par Mordell-Weil-Néron)}, donc l'hom
\[
(18)\quad \operatorname{Pic}(S)\longrightarrow\operatorname{Pic}(S)^\wedge\hookrightarrow H^2(S,T)
\]
est \uline{injectif}.

Sans nous soucier de cette condition, \deleted{\unsure{pour}} considérons le cas général --- je dis que la classe \deleted{\unsure{de}} $c$ \inserted{(14)}, est dans l'image de (18), de façon précise que c'est l'image de l'élément
\[
(19)\quad \unsure{\gamma_i|q_i}\in\operatorname{Pic}(S)
\]
classe du faisceau conormal (ou normal, le diable sait) de $X$ le long de $\unsure{q|g}_i$. Principe d'une vérification : on prend le complété formel de $X$ le long de $\unsure{q|g}_i(S)$ \unsure{et} \deleted{\unsure{le}} \unsure{qu'on ouvre} --- on enlève $\unsure{q|g}_i(S)$, et on interprète la suite exacte (13) comme la suite exacte d'homotopie de ce \deleted{schéma} topos \unsure{inventé|monté}, au dessus de \deleted{\unsure{Spec}} $S$. On a donc à prouver une histoire d'ombres\dots
\margin{80}
